\documentclass{article}
\usepackage{amsmath, amssymb, amsthm}
\usepackage{hyperref}
\usepackage{geometry}
\usepackage{titlesec}
\geometry{margin=1in}

\title{Erd\H{o}s Problems --- Verifiable\\[0.5em]\large A Collection of 7 Problems}
\date{Generated: 2026-06-09}
\author{Source: \url{https://www.erdosproblems.com}}

\begin{document}
\maketitle
\tableofcontents
\newpage


%% ============================================================
\section{Problem \#7}
\label{prob:7}

\noindent\textbf{Status:} VERIFIABLE \quad|\quad \textbf{Prize:} none \quad|\quad \textbf{Updated:} 2025-08-31

\noindent\textbf{Tags:} number theory, covering systems

\medskip

\noindent\textbf{Statement:}

Is there a distinct covering system all of whose moduli are odd?




Asked by Erd\H{o}s and Selfridge (sometimes also with Schinzel). They also asked whether there can be a covering system such that all the moduli are odd and squarefree. The answer to this stronger question is no, proved by Balister, Bollob\'{a}s, Morris, Sahasrabudhe, and Tiba \cite{BBMST22}. Hough and Nielsen \cite{HoNi19} proved that at least one modulus must be divisible by either $2$ or $3$. A simpler proof of this fact was provided by Balister, Bollob\'{a}s, Morris, Sahasrabudhe, and Tiba \cite{BBMST22}, who also prove that if an odd covering system exists then the least common multiple of its moduli must be divisible by $9$ or $15$. Selfridge has shown (as reported in \cite{Sc67}) that such a covering system exists if a covering system exists with moduli $n_1,\ldots,n_k$ such that no $n_i$ divides any other $n_j$ (but the latter has been shown not to exist, see [586] ). Filaseta, Ford, and Konyagin \cite{FFK00} report that Erd\H{o}s, 'convinced that an odd covering does exist, offered \$25 for a proof that no odd covering exists; Selfridge, convinced (at that point) that no odd covering exists, offered \$300 for the first explicit example...no award was promised to someone who gave a non-constructive proof that an odd covering of the integers exists...Selfridge (private communication) has informed us that he is now increasing his award to \$2000.'




References

[BBMST22] Balister, Paul and Bollob\'{a}s, B\'{e}la and Morris, Robert and Sahasrabudhe, Julian and Tiba, Marius, On the Erd\H{o}s covering problem: the density of the uncovered set . Invent. Math. (2022), 377-414.

[FFK00] Filaseta, M. and Ford, K. and Konyagin, S., On an irreducibility theorem of {A}. Schinzel associated with coverings of the integers . Illinois J. Math. (2000), 633--643.

[HoNi19] Hough, Robert D. and Nielsen, Pace P., Covering systems with restricted divisibility . Duke Math. J. (2019), 3261-3295.

[Sc67] Schinzel, A., Reducibility of polynomials and covering systems of congruences . Acta Arith. (1967/68), 91-101.

\noindent\rule{\linewidth}{0.4pt}


%% ============================================================
\section{Problem \#307}
\label{prob:307}

\noindent\textbf{Status:} VERIFIABLE \quad|\quad \textbf{Prize:} none \quad|\quad \textbf{Updated:} 2025-09-09

\noindent\textbf{Tags:} number theory, unit fractions

\medskip

\noindent\textbf{Statement:}

Are there two finite sets of primes $P,Q$ such that\[1=\left(\sum_{p\in P}\frac{1}{p}\right)\left(\sum_{q\in Q}\frac{1}{q}\right)?\]




Asked by Barbeau \cite{Ba76}. Can this be done if we drop the requirement that all $p\in P$ are prime and just ask for them to be relatively coprime, and similarly for $Q$? Cambie has found several examples when this weakened version is true. For example,\[1=\left(1+\frac{1}{5}\right)\left(\frac{1}{2}+\frac{1}{3}\right)\]and\[1=\left(1+\frac{1}{41}\right)\left(\frac{1}{2}+\frac{1}{3}+\frac{1}{7}\right).\]There are no examples known of the weakened coprime version if we insist that $1\not\in P\cup Q$. It is easy to see that, if $P$ and $Q$ are sets of primes, then $P$ and $Q$ are disjoint, and $\sum_{p\in P\cup Q}\frac{1}{p}\geq 2$, whence $\lvert P\cup Q\rvert \geq 60$.




References

[Ba76] Barbeau, E. J., Computer challenge corner: Problem 477: A brute force program . J. Rec. Math. (1976).

\noindent\rule{\linewidth}{0.4pt}


%% ============================================================
\section{Problem \#364}
\label{prob:364}

\noindent\textbf{Status:} VERIFIABLE \quad|\quad \textbf{Prize:} none \quad|\quad \textbf{Updated:} 2025-08-31

\noindent\textbf{Tags:} number theory

\medskip

\noindent\textbf{Statement:}

Are there any triples of consecutive positive integers all of which are powerful (i.e. if $p\mid n$ then $p^2\mid n$)?




Erd\H{o}s originally asked Mahler whether there are infinitely many pairs of consecutive powerful numbers, but Mahler immediately observed that the answer is yes from the infinitely many solutions to the Pell equation $x^2=8y^2+1$. This conjecture was also made by Mollin and Walsh \cite{MoWa86}. Erd\H{o}s \cite{Er76d} believed the answer to this question is no, and in fact if $n_k$ is the $k$th powerful number then\[n_{k+2}-n_k &gt; n_k^c\]for some constant $c&gt;0$. The abc conjecture implies there are only finitely many such triples. It is trivial that there are no quadruples of consecutive powerful numbers since one must be $2\pmod{4}$. Chan \cite{Ch25} has shown there are no triples $n-1,n,n+1$ of powerful numbers with $n$ a cube and $n-1,n+1$ both of the shape $p^3x^2$ for some prime $p$ and integer $x$. She \cite{Sh25} has proved there are no triples with $n$ a cube and both $n-1,n+1$ of the shape $p^2x^3$ with $p$ a prime. By OEIS A076445 there are no such $n$ for $n&lt;7.38\times 10^{28}$. See also [137] , [365] , and [938] .




References

[Ch25] Chan, Tsz Ho, A note on three consecutive powerful numbers . Integers (2025), Paper No. A7, 7.

[Er76d] Erd\H{o}s, P., Problems and results on number theoretic properties of consecutive integers and related questions . Proceedings of the Fifth Manitoba Conference on Numerical Mathematics (Univ. Manitoba, Winnipeg, Man., 1975) (1976), 25-44.

[MoWa86] Mollin, R. A. and Walsh, P. G., On powerful numbers . Internat. J. Math. Math. Sci. (1986), 801--806.

[Sh25] She, Jialai, Nonexistence of consecutive powerful triplets around cubes with prime-square factors . Integers (2025), Paper No. A103, 9.

\noindent\rule{\linewidth}{0.4pt}


%% ============================================================
\section{Problem \#366}
\label{prob:366}

\noindent\textbf{Status:} VERIFIABLE \quad|\quad \textbf{Prize:} none \quad|\quad \textbf{Updated:} 2025-08-31

\noindent\textbf{Tags:} number theory

\medskip

\noindent\textbf{Statement:}

Are there any $2$-full $n$ such that $n+1$ is $3$-full? That is, if $p\mid n$ then $p^2\mid n$ and if $p\mid n+1$ then $p^3\mid n+1$.




Erd\H{o}s originally asked Mahler whether there are infinitely many pairs of consecutive powerful numbers, but Mahler immediately observed that the answer is yes from the infinitely many solutions to the Pell equation $x^2=8y^2+1$. Note that $8$ is 3-full and $9$ is $2$-full. Erd\H{o}s and Graham asked if this is the only pair of such consecutive integers. Stephan has observed that $12167=23^3$ and $12168=2^33^213^2$ (a pair already known to Golomb \cite{Go70}) is another example, but (by OEIS A060355 ) there are no other examples for $n&lt;10^{22}$. In \cite{Er76d} Erd\H{o}s asks the weaker question of whether there are any consecutive pairs of $3$-full integers (which is also discussed in problem B16 of Guy's collection \cite{Gu04}). Turturean notes in the comments that the ABC conjecture implies there are only finite ly many such $n$.




References

[Er76d] Erd\H{o}s, P., Problems and results on number theoretic properties of consecutive integers and related questions . Proceedings of the Fifth Manitoba Conference on Numerical Mathematics (Univ. Manitoba, Winnipeg, Man., 1975) (1976), 25-44.

[Go70] Golomb, S. W., Powerful numbers . Amer. Math. Monthly (1970), 848-855.

[Gu04] Guy, Richard K., Unsolved problems in number theory . (2004), xviii+437.

\noindent\rule{\linewidth}{0.4pt}


%% ============================================================
\section{Problem \#647}
\label{prob:647}

\noindent\textbf{Status:} VERIFIABLE \quad|\quad \textbf{Prize:} £25 \quad|\quad \textbf{Updated:} 2025-08-31

\noindent\textbf{Tags:} number theory

\medskip

\noindent\textbf{Statement:}

Let $\tau(n)$ count the number of divisors of $n$. Is there some $n&#62;24$ such that\[\max_{m&#60;n}(m+\tau(m))\leq n+2?\]




A problem of Erd\H{o}s and Selfridge. This is true for $n=24$. The $n+2$ is best possible here since\[\max(\tau(n-1)+n-1,\tau(n-2)+n-2)\geq n+2.\]In \cite{Er79} Erd\H{o}s says 'it is extremely doubtful' that there are infinitely many such $n$, and in fact suggets that\[\lim_{n\to \infty}\max_{m&lt;n}(\tau(m)+m-n)=\infty.\]In \cite{Er79d} Erd\H{o}s says it 'seems certain' that for every $k$ there are infinitely many $n$ for which\[\max_{n-k&lt;m&lt;n}(m+\tau(m))\leq n+2,\]but 'this is hopeless with our present methods', although it follows from Schinzel's Hypothesis H . In \cite{Er92e} Erd\H{o}s offered £25 for an example of such an $n&gt;24$. He wrote 'I am being rather stingy but we old people are stingy.' (This has been converted to \$44 using approximate 1992 exchange rates.) Tao has observed in the comments that, since $\tau(m)$ is similar to $2^{\omega(m)}$, this problem is similar to (but slightly weaker than) the first part of [679] , but much stronger than [413] or [248] . See also [413] .




References

[Er79] Erd\H{o}s, Paul, Some unconventional problems in number theory . Math. Mag. (1979), 67-70.

[Er79d] Erd\H{o}s, P., Some unconventional problems in number theory . Acta Math. Acad. Sci. Hungar. (1979), 71-80.

[Er92e] Erd\H{o}s, P\'{a}l, Some Unsolved problems in Geometry, Number Theory and Combinatorics . Eureka (1992), 44-48.

\noindent\rule{\linewidth}{0.4pt}


%% ============================================================
\section{Problem \#672}
\label{prob:672}

\noindent\textbf{Status:} VERIFIABLE \quad|\quad \textbf{Prize:} none \quad|\quad \textbf{Updated:} 2025-08-31

\noindent\textbf{Tags:} number theory

\medskip

\noindent\textbf{Statement:}

Can the product of an arithmetic progression of positive integers $n,n+d,\ldots,n+(k-1)d$ of length $k\geq 4$ (with $(n,d)=1$) be a perfect power?




Erd\H{o}s believed not. Erd\H{o}s and Selfridge \cite{ErSe75} proved that the product of consecutive integers is never a perfect power. The theory of Pell equations implies that there are infinitely many pairs $n,d$ with $(n,d)=1$ such that $n(n+d)(n+2d)$ is a square. Considering the question of whether the product of an arithmetic progression of length $k$ can be equal to an $\ell$th power: {UL} {LI}Euler proved this is impossible when $k=4$ and $\ell=2$,{/LI} {LI}Obl\'{a}th \cite{Ob51} proved this is impossible when $(k,l)=(5,2),(3,3),(3,4),(3,5)$.{/LI} {LI}Marszalek \cite{Ma85} proved that this is only possible for $k\ll_d 1$, where $d$ is the common difference of the arithmetic progression.{/LI} {LI}Gy\"{o}ry, Hajdu, and Saradha \cite{GHS04} proved this is impossible for $4\leq k\leq 5$. {LI}Bennett, Bruin, Gy\"{o}ry, and Hajdu \cite{BBGH06} proved this is impossible for $4\leq k\leq 11$, and also impossible for $k$ sufficiently large depending only on the number of prime divisors of $d$.{/LI} {LI} Gy\"{o}ry, Hajdu, and Pint\'{e}r \cite{GHP09} have proved this is impossible for $4\leq k\leq 34$.{/LI} {LI} Bennett and Siksek \cite{BeSi20} have proved that this is impossible for all $k$ sufficiently large and $\ell &#62; e^{10^k}$ prime.{/LI} {/UL} Jakob F\"{u}hrer has observed this is possible for integers in general, for example $(-6)\cdot(-1)\cdot 4\cdot 9=6^3$.




References

[BBGH06] Bennett, M. A. and Bruin, N. and Gy\H ory, K. and Hajdu, L., Powers from products of consecutive terms in arithmetic progression . Proc. London Math. Soc. (3) (2006), 273--306.

[BeSi20] Bennett, Michael A. and Siksek, Samir, A conjecture of {E}rd\H os, supersingular primes and short character sums . Ann. of Math. (2) (2020), 355--392.

[ErSe75] Erd\H{o}s, P. and Selfridge, J. L., The product of consecutive integers is never a power . Illinois J. Math. (1975), 292-301.

[GHP09] Gy\H ory, K. and Hajdu, L. and Pint\'er, \'A., Perfect powers from products of consecutive terms in arithmetic progression . Compos. Math. (2009), 845--864.

[GHS04] Gy\H ory, K. and Hajdu, L. and Saradha, N., On the {D}iophantine equation {$n(n+d)\cdots(n+(k-1)d)=by^l$} . Canad. Math. Bull. (2004), 373--388.

[Ma85] Marsza\l ek, R., On the product of consecutive elements of an arithmetic progression . Monatsh. Math. (1985), 215-222.

[Ob51] Oblath, Richard, Eine Bemerkung \"{u}ber Produkte aufeinander folgender Zahlen . J. Indian Math. Soc. (N.S.) (1951), 135-139.

\noindent\rule{\linewidth}{0.4pt}


%% ============================================================
\section{Problem \#835}
\label{prob:835}

\noindent\textbf{Status:} VERIFIABLE \quad|\quad \textbf{Prize:} none \quad|\quad \textbf{Updated:} 2025-08-31

\noindent\textbf{Tags:} graph theory, hypergraphs

\medskip

\noindent\textbf{Statement:}

Does there exist a $k&#62;2$ such that the $k$-sized subsets of $\{1,\ldots,2k\}$ can be coloured with $k+1$ colours such that for every $A\subset \{1,\ldots,2k\}$ with $\lvert A\rvert=k+1$ all $k+1$ colours appear among the $k$-sized subsets of $A$?




A problem of Erd\H{o}s and Rosenfeld. This is trivially possible for $k=2$. They were not sure about $k=6$. This is equivalent to asking whether there exists $k&gt;2$ such that the chromatic number of the Johnson graph $J(2k,k)$ is $k+1$ (it is always at least $k+1$ and at most $2k$). The chromatic numbers listed at this website show that this is false for $3\leq k\leq 8$. Ma and Tang have proved that the chromatic number of $J(2k,k)$ is $&gt;k+1$ for all $k&gt;2$ not of the form $p-1$ for prime $p$.

\noindent\rule{\linewidth}{0.4pt}

\end{document}
